\section{Harness and Run Protocol}
\label{sec:harness}

\paragraph{No automation framework.} A framework that drives the browser
(Puppeteer, Playwright, a debugging protocol) attaches code and threads to
the process being measured. The harness instead launches the browser binary
like a user would, with the page URL carrying the configuration, and the
page posts its results to a local server. The server also sends the COOP
and COEP headers that shared memory needs. Nothing is attached to the
browser while it measures.

\paragraph{A browser run.} The launcher creates an empty temporary profile,
starts the browser headless in its own process group, inside the run's
scope and mask, and waits for the result file. The page waits a fixed start
delay, creates the coordinating worker (the main thread cannot block),
which instantiates the module, creates the thread pool and runs the
measurement, and posts the rows and the page's metadata. The launcher then
stops the scope, signals the whole process group so that no renderer
survives into the next run, and deletes the profile. Background work is
turned off symmetrically in both browsers (updates, telemetry, safe
browsing, remote settings, prefetching; background throttling in Chrome).
A run ends as done, timed out or exited without a result; a failed run is
recorded, not retried silently, and the campaign reports a failure rate.

\paragraph{A native run} goes through a wrapper that starts the runner in
the same kind of scope with the same mask and writes the same pair of files
(the rows and the context).

\begin{figure}[t]
\centering
\resizebox{\linewidth}{!}{%
\begin{tikzpicture}[font=\footnotesize, node distance=0pt,
  ph/.style={draw, minimum height=8mm, inner sep=3pt, align=center}]
  \node[ph, fill=gray!10, minimum width=22mm] (launch) {launch\\(+ 3 s delay)};
  \node[ph, fill=gray!10, minimum width=20mm, right=of launch] (pool) {input,\\thread pool};
  \node[ph, fill=gray!25, minimum width=24mm, right=of pool] (cal) {fork-join\\calibration};
  \node[ph, fill=gray!25, minimum width=26mm, right=of cal] (warm) {warm-up\\$W$ iterations};
  \node[ph, fill=native!25, minimum width=44mm, right=of warm] (meas) {10 measured iterations\\run value = their median};
  \node[ph, fill=gray!10, minimum width=20mm, right=of meas] (post) {write or\\post rows};
  \draw[decorate] ($(warm.south west)+(0,-3mm)$) -- node[below, font=\scriptsize, yshift=-1mm] {one iteration: \code{prepare} (untimed), \code{run} (timed), \code{checksum} (untimed, one thread)} ($(meas.south east)+(0,-3mm)$);
  \draw ($(warm.south west)+(0,-3mm)$) -- ++(0,1.5mm);
  \draw ($(meas.south east)+(0,-3mm)$) -- ++(0,1.5mm);
  \node[font=\scriptsize, above=1mm of pool.north east, anchor=south] {memory sample 1};
  \node[font=\scriptsize, above=1mm of meas.north west, anchor=south] {sample 2};
  \node[font=\scriptsize, above=1mm of meas.north east, anchor=south] {sample 3};
\end{tikzpicture}}
\caption{The layout of one run, the same in the three environments. Only
the measured iterations enter the analysis; the calibration is reported
separately and nothing is subtracted from the times.}
\Description{A timeline of a run: launch, input and thread pool, fork-join calibration, warm-up iterations, ten measured iterations, and writing the rows.}
\label{fig:run}
\end{figure}

\paragraph{The layout of a run} is shared code, the same in the three
environments (Figure~\ref{fig:run}). The input is generated and the pool of
$N$ threads is created first. Then the \emph{fork-join calibration} times
empty parallel regions, 20 samples each of two probes: an empty
\code{broadcast}, which every worker must run (a barrier), and an empty
parallel iterator (ordinary optional parallelism); each probe is timed warm
(back to back, workers spinning) and cold (after a pause at least as long
as the workload's checksum, so the workers are asleep, as they are before a
measured iteration). Then come $W$ warm-up iterations and 10 measured
ones. In each iteration only \code{run} is between the two clock reads.
Every iteration is called from outside the pool, natively from the main
thread and in WASM from the coordinating worker, as an application calls a
parallel library. The memory size (the native resident set, the WASM linear
memory) is recorded before the warm-up, before the first measured iteration
and after the last; a WASM run whose memory grows during the measured phase
fails. The clock is \code{performance.now()} in the browsers, with a
resolution measured in every run: 5~\textmu s in Chrome, 20~\textmu s in
Firefox.

\paragraph{What a run writes.} One CSV with a row per iteration (phase,
index, nanoseconds, checksum, and the settings of the run) and one JSON
with the context of Section~\ref{sec:platform}. The run identifier is not
stored in the rows; the analysis joins it, so a file cannot be filed under
the wrong run.

\paragraph{The orchestrator.} A campaign is a JSON specification holding
every parameter: environments, workloads with their input, $k$ and warm-up,
thread counts, repetitions, image size, iterations, calibration samples,
masks, start delay and timeout. Unknown keys are rejected. The orchestrator
builds every combination, shuffles the whole list once with the recorded
seed and runs it in that order, so that a slow drift of the machine
(temperature, background activity) is spread over the environments instead
of lining up with one. Before the plan, each browser is launched once and
discarded, because the first launch on a machine is slow. A campaign that
is stopped resumes: finished runs are skipped and the order is the same. It
refuses to continue with another specification, commit, artifact digest or
browser version, so one campaign never mixes builds. During a campaign
nothing else is started on the machine.

\paragraph{The session of a build.} Measuring a build is a sequence of
numbered, resumable stages, each writing its log:
\begin{enumerate}[leftmargin=2em]
\item \emph{Preflight}: the host state as found, the environment set up and
checked, the tests, a tiny copy of every campaign to exercise the pipeline
end to end, and the per-core check of Section~\ref{sec:core}.
\item \emph{Build}: the tests, two builds from scratch, the digests, the
tag.
\item \emph{Gate}: a short campaign on the cells whose code changed, with
criteria fixed in advance (Section~\ref{sec:gate}).
\item \emph{Main campaign}, then the exploratory ones, in an order chosen so
that what is dropped for lack of time is dropped from the end.
\item \emph{Results}: every table and figure, then a check that
regenerating them changes nothing.
\end{enumerate}

\subsection{The gate}
\label{sec:gate}

The gate ran the cells whose code changed (both histograms at $k=1$ and
256, the sequential workloads) and the filter as a control, at $N=1$ and 8
in the three environments, 3 runs each at the real size: 135 runs. It
passes when, in every cell, the checksums agree across environments and
thread counts, every timed sample is at least 100 clock ticks, and the
median of the first three measured iterations is within 5\% of the median
of the last three. Three runs only catch a gross failure, such as a warm-up
that no longer covers a tier change. How a failure is read was fixed
beforehand too: the gate is run once more, whole, with another seed, and
the same cell failing twice stops the session for a recorded decision.

Both gates ran without a failed or excluded run and passed the first two
criteria; 44 of 45 cells passed the third, and the same cell failed both
times: Firefox, local histogram, $k=1$, $N=8$, with ratios of 1.052 and
1.077. In every one of its six runs, one or two of the first three measured
iterations take 21 to 22~ms and the others 19.7~ms; the 71 warm-up
iterations before them show nothing of the kind. The run value is a median
of ten iterations, which one or two slow ones do not move: it is 1.000 to
1.0015 times the value taken from the last seven iterations alone. The
decision, recorded before the main campaign started: the campaigns go ahead
with every parameter unchanged, the gate stays recorded as failed on one
criterion in one cell, and the cell is declared with its measured effect of
about 0.1\%. Changing the measurement after seeing data (dropping the first
iterations, lengthening the warm-up) was the alternative rejected. The tier
campaign later showed that these slow iterations come with the optimizing
compiler's code (27 of 30 runs as shipped and with the optimizing compiler
only, 1 of 30 with the baseline compiler only); what they are is not known.
